\FifthArticle{10}

Es bleibt nur noch nachzuweisen: Liegt der eine Wurzel $=p$\SourceLeaf{090} darstellende Punkt $P$ zugleich auf zwei oder mehr Linien $T=0$, so ist das Quadrat von $x-p$ oder die der Anzahl jener zusammentreffenden Linien entsprechende höhere Potenz als Faktor in $X$ enthalten. Der Beweis dafür beruht auf den folgenden Sätzen.

Man führe anstelle der unbestimmten Größe $x$ eine andere $z$ ein, indem man $x=z+p$ setzt. Durch diese Substitution gehe $X$ in $Z$ über. Dann ist $Z$ eine Funktion von $z$ vom gleichen Grad wie $X$ in $x$, deren konstantes Glied aber fehlt. Ordnet man sie nach aufsteigenden Potenzen von $z$, so sei das niedrigste nicht verschwindende Glied
\[
=Kz^m\quad\mbox{und}\quad Z=Kz^m(1+\zeta)
\]
wobei $\zeta$ die Form $Lz+L'zz+L''z^3+\mbox{usw.}+\frac{1}{K}z^{n-m}$ hat. Schließlich setze man
\[
z=s(\cos\psi+i\sin\psi)
\]

Der reelle und der imaginäre Bestandteil von $z$ drücken als rechtwinklige Koordinaten die Lage jedes unbestimmten Punktes der Ebene aus, und die Größen $s$, $\psi$ als Polarkoordinaten, jeweils bezüglich des Punktes $P$. Das entspricht genau der Weise, in der die Bestandteile von $x$ und die Größen $r$, $\varphi$ die Lage bezüglich des ursprünglichen Koordinatenursprungs bezeichnen. Verbindet man einen bestimmten Wert von $s$ mit allen Werten von $\psi$ über einen Bereich von $360^0$, so werden dadurch die Punkte einer Kreislinie dargestellt, deren Mittelpunkt $P$ und deren Radius $=s$ ist.

Setzt man nun $K=k(\cos\kappa+i\sin\kappa)$ und folglich
\[
Kz^m=ks^m(\cos(m\psi+\kappa)+i\sin(m\psi+\kappa))
\]
so darf für ein unendlich kleines $s$ die Größe $\zeta$, die mindestens von derselben Ordnung wie $s$ ist, gegenüber $1$ vernachlässigt und somit
\[
T=ks^m\cos(m\psi+\kappa)
\]
gesetzt werden. Daraus folgt, dass, während $\psi$ um $360^0$ wächst, das Vorzeichen von $T$ auf $m$ Stücken des Kreisumfangs positiv und auf ebenso vielen mit ihnen abwechselnden negativ ist, oder dass $T$ an $2m$ Punkten $=0$ wird, nämlich für $\psi=\frac{1}{m}(\kappa-90^0)$, $\frac{1}{m}(\kappa+90^0)$, $\frac{1}{m}(\kappa+270^0)$ usw. Von $P$\SourceLeaf{091} gehen demnach insgesamt $2m$ Linien aus, in denen $T=0$ ist. Verbindet man sie aber paarweise so, dass jede, bei der mit wachsendem $\psi$ das Vorzeichen von $-$ in $+$ übergeht, zusammen mit der nächstfolgenden, bei der der entgegengesetzte Übergang stattfindet, als Begrenzungslinie eines Flächenteils mit positivem $T$ betrachtet wird, so treffen in $P$ insgesamt $m$ solche Begrenzungslinien zusammen.

Andererseits ist klar: So wie $Z$ identisch durch $z^m$ und durch keine höhere Potenz von $z$ teilbar ist, enthält $X$ den Faktor $(x-p)^m$, aber keine höhere Potenz von $x-p$. Wenn also $p$ irgendeine Wurzel der Gleichung $X=0$ bezeichnet, ist der Exponent der höchsten Potenz von $x-p$, durch die $X$ teilbar ist, immer gleich der Anzahl der \emph{in} $P$ zusammentreffenden Begrenzungslinien für Flächen mit positivem $T$, oder, was dasselbe ist, der Anzahl solcher \emph{an} $P$ zusammentreffender Flächen.

Übrigens lässt sich der Beweisführung leicht eine Form geben, die von der Einbeziehung unendlich kleiner Größen völlig unabhängig und der Schlussreihe in den Artikeln 3 und 4 ganz ähnlich ist. Es lässt sich nämlich ein Wert von $s$ nachweisen, für den ebenso wie für jeden kleineren der gesamte Zyklus aller Werte von $\psi$ dieselbe abwechselnde Folge von $m$ Stücken mit positivem $T$ und ebenso vielen mit negativem Wert bietet. Diese Eigenschaft hat die positive Wurzel der Gleichung
\[
0=m\surd\frac{1}{2}-(m+1)ls-(m+2)l'ss-(m+3)l''s^3-\mbox{usw.}
\]
wobei $l$, $l'$, $l''$ usw. die positiven Quadratwurzeln aus den Normen der komplexen Größen $L$, $L'$, $L''$ usw. bedeuten, oder wobei
\[
\begin{array}{rl}
L&=l\,(\cos\lambda+i\sin\lambda)\\
L'&=l'(\cos\lambda'+i\sin\lambda')\\
L''&=l''(\cos\lambda''+i\sin\lambda'')\ \mbox{usw.}
\end{array}
\]
gesetzt ist. Ich glaube jedoch, die sehr leichte Entwicklung dieses Satzes hier übergehen zu können.

Schließlich sei noch bemerkt, dass bei der Beweisführung in der Abhandlung von 1799 die Betrachtung \emph{zweier} Systeme von Linien erforderlich war: Das eine enthielt die Linien, in denen $T=0$, das andere diejenigen, in denen $U=0$ ist. In unserem jetzigen Verfahren hat dagegen die Betrachtung eines Systems\SourceLeaf{092} ausgereicht. Ich habe dafür das System der Begrenzungslinien der Flächenteile mit positivem $T$ gewählt; ebenso gut hätte aber die Betrachtung der Begrenzungslinien der Flächen mit positivem (oder negativem) $U$ demselben Zweck dienen können.

